% =============================================================================
% Results
% =============================================================================
\section{Results}\label{sec:results}

\subsection{Eccentricity sweep within the Ritz model}

The principal result is the dependence of $\kappa$ on eccentricity
\emph{within the same model}.  Using the oblate Ritz forward model
(Section~\ref{sec:perturbation}) and the scaled Jacobian built from flexural
modes $n=2,\ldots,6$, we sweep the aspect ratio $\zeta=c/a$ from near-unity
(almost spherical) to the canonical value $\zeta=0.667$
($\varepsilon=0.745$).  At the near-spherical limit
($\varepsilon \to 0$) the condition number rises to the floor value
\begin{equation}\label{eq:kappa_floor}
  \kappa_\mathrm{floor} \approx 269
  \qquad
  (\text{evaluated at } \varepsilon = 0.01,\;
  \zeta = 0.99995,\;\text{five modes}),
\end{equation}
representing the curvature-channel contribution identified in
Proposition~\ref{prop:eccentricity}.  At the canonical oblate operating point
the condition number drops to
\begin{equation}
  \kappa_\mathrm{oblate} = 69.4,
\end{equation}
an improvement of approximately one order of magnitude achieved purely through
geometric eccentricity.  This Ritz-to-Ritz comparison demonstrates that it is
\emph{geometry}, not model choice, that drives the identifiability transition.

\begin{remark}[Equivalent-sphere comparison as pedagogical aside]
For context, the equivalent-sphere model of
Section~\ref{sec:theory} gives $\kappa_\mathrm{sphere}=1.37\times10^{10}$
at the same parameter set, reflecting the algebraic rank deficiency of
Proposition~\ref{prop:chain_rule} rather than the eccentricity of the shell.
The eight-order-of-magnitude gap between $\kappa_\mathrm{sphere}$ and
$\kappa_\mathrm{oblate}$ is dramatic, but it conflates two effects ---
model architecture and physical geometry --- and should not be read as
``oblateness reduces the condition number by $10^8$''.
The Ritz eccentricity sweep above isolates the geometric effect alone.
\end{remark}

\begin{figure}[t]
  \centering
  \includegraphics[width=\linewidth]{fig_sphere_vs_oblate}
  \caption{Frequency spectra comparison between the near-spherical and
  canonical oblate operating points for flexural modes $n=2,\ldots,6$,
  computed within the oblate \emph{Ritz} model.  The oblate case produces a
  well-separated spectrum; the near-spherical case shows near-degenerate
  clustering.}
  \label{fig:sphere_vs_oblate}
\end{figure}

\subsection{Singular-value structure}

The singular value decomposition provides the clearest picture.  In the
equivalent-sphere model, one singular value collapses to a value near machine
zero, leaving one parameter combination effectively unobservable.  In the
oblate Ritz model at $\varepsilon=0.745$, all three singular values remain
finite and well separated from zero.  Even at the near-spherical Ritz limit
($\varepsilon\to0$), all three singular values remain non-zero --- the
curvature-channel floor identified in Proposition~\ref{prop:eccentricity}
prevents complete rank collapse --- though the smallest value is markedly
reduced, yielding $\kappa_\mathrm{floor}\approx269$.

This observation is more informative than the condition number alone.  A large
$\kappa$ could arise from one weak direction in an otherwise healthy problem; in
the equivalent-sphere case, the singular values show something stronger: the
Jacobian is algebraically rank deficient by construction
(Proposition~\ref{prop:chain_rule}).  At the canonical oblate Ritz operating
point, by contrast, the inverse map is not perfect but it is usable.

\begin{figure}[t]
  \centering
  \includegraphics[width=0.85\linewidth]{fig_singular_values}
  \caption{Singular value spectrum of the scaled Jacobian at the
  near-spherical (\emph{Ritz} model, $\varepsilon\approx0$) and
  canonical oblate (\emph{Ritz} model, $\varepsilon=0.745$) operating points,
  with the \emph{equivalent-sphere} model shown for comparison.
  In the equivalent-sphere case, the smallest singular value
  collapses to near-zero, rendering the Jacobian numerically rank deficient.
  In the oblate Ritz case, all three singular values remain well separated from
  zero.}
  \label{fig:singular_values}
\end{figure}

\subsection{Physical interpretation}

Two models are compared in this paper.  The equivalent-sphere model
(Section~\ref{sec:theory}) depends on $(a,c)$ only through
$R_\mathrm{eq}=(a^2c)^{1/3}$; its Jacobian is algebraically rank deficient
by Proposition~\ref{prop:chain_rule}.  The oblate Ritz model
(Section~\ref{sec:perturbation}) integrates stiffness and mass over the true
spheroidal surface, so each mode samples curvature differently.  All
condition-number results in this section --- including the eccentricity sweep,
the Cram\'er--Rao bounds, and the round-trip inversions ---
use the Ritz model unless explicitly stated otherwise.

Within the Ritz model, the modal balance of bending
($K_\mathrm{bend}^{(n)} \propto R^{-4}$), membrane
($K_\mathrm{memb}^{(n)} \propto R^{-2}$), prestress
($K_\mathrm{pre}^{(n)} \propto R^{-1}$), and added mass
($m_\mathrm{eff}^{(n)} = \rho_w h + \rho_f R/n$) varies
sufficiently with mode number to separate the sensitivity directions away from
the spherical limit.  The key mechanism is the mode-dependent weighting of the
meridional curvature integral: different $P_n(\eta)$ patterns sample the
equatorial and polar regions with different emphasis, producing $n$-dependent
corrections to $\partial f_n/\partial a$ and $\partial f_n/\partial c$ that
break the proportionality forced by $R_\mathrm{eq}$-only dependence.

\subsection{Round-trip verification}

As a sanity check on the numerical Jacobian and solver implementation, we
perform round-trip inversion: synthetic frequencies are generated by the Ritz
forward model and the solver is initialised within the local basin of
attraction.  The recovered parameters satisfy
\begin{equation}
  \frac{|\hat a-a|}{a},
  \frac{|\hat c-c|}{c},
  \frac{|\hat E-E|}{E}
  < 10^{-12}.
\end{equation}
This machine-precision recovery confirms that the forward model is locally
invertible and that the numerical Jacobian is consistent, but it should not be
over-interpreted.  Because the same model generates and inverts the data, the
test is tautological with respect to model-form error: it verifies internal
consistency, not the ability to recover parameters from real measurements.
The round-trip result is therefore a necessary verification, not a key finding.

The distinction between local and global behaviour is important.  With only
three modes, the inverse problem is exactly determined and the solver can
encounter local minima or alternative low-residual basins.  With five modes
($n=2,\ldots,6$), the system becomes overdetermined and the same inversion is
globally robust across the tested perturbations of the initial guess.  The
results therefore support a clear experimental design rule: three modes may be
enough to tempt the optimiser, but five modes are what persuade it to behave.

\subsection{Cram\'er--Rao bounds}

Table~\ref{tab:crlb} summarises the Cram\'er--Rao lower bounds for the canonical
oblate case at $1\%$ frequency noise.  The smallest uncertainty is obtained for
$c$, whereas $a$ and $E$ are comparably less well constrained.

\begin{table}[t]
  \centering
  \caption{Cram\'er--Rao lower bounds for the canonical oblate case at
  $1\%$ fractional frequency noise, using five flexural modes
  ($n=2,\ldots,6$).  Relative bounds are reported as one
  standard deviation.}
  \label{tab:crlb}
  \begin{tabular}{lccc}
    \toprule
    Parameter & Nominal value & Absolute bound & Relative bound \\
    \midrule
    $a$ & \SI{0.18}{m}    & $\pm\,\SI{0.040}{m}$   & $\pm\,22.5\%$ \\
    $c$ & \SI{0.12}{m}    & $\pm\,\SI{0.007}{m}$   & $\pm\,6.0\%$ \\
    $E$ & \SI{0.1}{MPa}   & $\pm\,\SI{0.031}{MPa}$ & $\pm\,30.8\%$ \\
    \bottomrule
  \end{tabular}
\end{table}

The $\pm30.8\%$ bound on Young's modulus merits contextualisation.  In
quasi-static elastography of abdominal soft tissue, reported measurement
uncertainties typically range from $\pm20\%$ to $\pm50\%$, depending on organ,
technique, and boundary conditions~\cite{Doyley2012}.  Shear-wave
elastography of the liver, for instance, reports inter-observer coefficients of
variation of $15$--$25\%$~\cite{Barr2015}, while in vivo mechanical testing of
the abdominal wall yields variability of $30$--$60\%$ across
subjects~\cite{Aernouts2012}.  The CRB of $\pm30.8\%$ therefore falls within
the range of existing clinical measurement precision: the spectral inverse
problem does not promise better resolution than elastography already achieves,
but neither does it demand unrealistic frequency measurement accuracy.

For the equivalent-sphere case, the Fisher matrix is singular and no finite
Cram\'er--Rao bound exists in the degenerate direction.  This is the statistical
counterpart of the singular-value result: the data simply do not contain enough
independent information at that operating point.

As expected from \eqref{eq:fisher}, the bounds scale linearly with noise level.
Halving the frequency noise from $2\%$ to $1\%$ halves the minimum-achievable
parameter standard deviations.  The uncertainty budget is therefore governed by
two levers: the intrinsic conditioning of the operating point and the
experimental precision of the measured resonances.

\begin{figure}[t]
  \centering
  \includegraphics[width=\linewidth]{fig_inversion_noise}
  \caption{Monte Carlo inversion error as a function of fractional frequency
  noise level.  Relative parameter recovery errors for $a$, $c$, and~$E$ are
  shown, confirming the linear scaling predicted by the Cram\'er--Rao bounds.
  Five-mode inversion remains robust across the tested noise range.}
  \label{fig:inversion_noise}
\end{figure}

\subsection{Mode-count robustness}\label{sec:mode_count}

Table~\ref{tab:mode_count} quantifies the dependence of key identifiability
metrics on the number of measured flexural modes, from the minimum
exactly-determined case ($\mathcal{N}=3$ modes) through the overdetermined
cases up to $\mathcal{N}=9$.  The condition number remains stable within
the range $\kappa\in[69,94]$ across all mode counts, confirming that the
rank restoration is a robust feature of the Ritz model geometry rather than
an artefact of a particular mode selection.

The Cram\'er--Rao bounds improve monotonically with mode count, as expected
from the increasing information content of additional measurements.  However,
the marginal improvement diminishes rapidly: the bound on $E$ drops from
$\pm36.2\%$ at three modes to $\pm30.8\%$ at five modes, and improves by less
than one percentage point thereafter.  Five modes therefore capture most of the
available information for the present three-parameter problem.

\begin{table}[t]
  \centering
  \caption{Identifiability metrics as a function of the number of
  \emph{measured} flexural modes $\mathcal{N}$, using modes
  $n=2,\ldots,\mathcal{N}+1$ in the Ritz model at the canonical operating
  point with $1\%$ frequency noise.  This is a measured-mode sensitivity
  analysis, not a Ritz basis convergence study: the Ritz trial-function
  order is held fixed while the number of observed modes varies.}
  \label{tab:mode_count}
  \begin{tabular}{ccccc}
    \toprule
    $\mathcal{N}$ & $\kappa$ & CRB($a$) & CRB($c$) & CRB($E$) \\
    \midrule
    3 &  73.0 & $\pm\,36.9\%$ & $\pm\,13.3\%$ & $\pm\,36.2\%$ \\
    5 &  69.4 & $\pm\,22.5\%$ & $\pm\,6.0\%$  & $\pm\,30.8\%$ \\
    7 &  80.6 & $\pm\,21.0\%$ & $\pm\,5.3\%$  & $\pm\,29.9\%$ \\
    9 &  93.6 & $\pm\,20.8\%$ & $\pm\,5.2\%$  & $\pm\,29.8\%$ \\
    \bottomrule
  \end{tabular}
\end{table}

\subsection{Condition-number map in parameter space}

Condition-number maps over $(a,\zeta_g)$ show a broad region of acceptable
conditioning across physiologically relevant oblate geometries.  The condition
number increases steadily as $\zeta_g \rightarrow 0$, that is, as the geometry
approaches a sphere.  This trend is physically and numerically consistent with
the singular spherical limit.  Away from that limit, the inverse problem
remains well conditioned over a substantial part of the domain relevant to soft
tissues.

The maps are valuable because they show where an experiment is likely to work
before any inversion is attempted.  They also suggest that identifiability is
not a single number attached to the model in general; it depends strongly on the
operating point.

\begin{figure}[t]
  \centering
  \includegraphics[width=0.85\linewidth]{fig_condition_map}
  \caption{Condition number $\kappa(\mathbf{J}_s)$ mapped over the
  $(a,\,\zeta_g)$ parameter plane using the \emph{Ritz} model.  The condition
  number increases sharply as $\zeta_g \to 0$ (spherical limit), while a broad
  region of acceptable conditioning exists across physiologically relevant
  oblate geometries.}
  \label{fig:condition_map}
\end{figure}

\subsection{Watermelon cross-application}

The same analysis applied to a ripe watermelon gives
\begin{equation}
  \kappa \approx 128.
\end{equation}
This value is somewhat worse than the canonical abdomen, but it is still well
within the range of practical inversion.  The result shows that the underlying
identifiability mechanism is not confined to one biological parameter set.

The watermelon case also helps separate geometry from scale.  The stiffness is
orders of magnitude larger than in soft tissue, yet the inverse problem remains
tractable.  The conditioning therefore reflects the structure of the
frequency--parameter map rather than a peculiarity of one very soft shell.

\subsection{Eccentricity dependence and power-law scaling}

The analytical predictions of Propositions~\ref{prop:chain_rule}
and~\ref{prop:eccentricity} are confirmed numerically.
Figure~\ref{fig:kappa_vs_eccentricity} shows the condition number
$\kappa(\mathbf{J}_s)$ as a function of eccentricity
$\varepsilon = \sqrt{1 - c^2/a^2}$, computed by sweeping the aspect ratio
$\zeta = c/a$ while holding all other parameters at their canonical values.
All computations use the oblate Ritz model.

As $\varepsilon \to 0$ (spherical limit), $\kappa$ rises steeply towards its
near-spherical floor of approximately $269$ (five modes, evaluated at
$\varepsilon = 0.01$, i.e.\ $\zeta = c/a = 0.99995$), consistent with
the curvature-channel contribution identified in
Proposition~\ref{prop:eccentricity}.  As
$\varepsilon$ increases (more oblate), $\kappa$ drops rapidly as the
mode-dependent curvature sampling of Proposition~\ref{prop:eccentricity}
restores full rank.  In the intermediate eccentricity regime (where the
shape-channel correction dominates the curvature-channel floor), the data are
well described by the power law
\begin{equation}
  \kappa \approx C\,\varepsilon^{-\alpha}.
\end{equation}

The theoretical analysis of Section~\ref{sec:perturbation} predicts
$\alpha=2$ from the $\varepsilon^2$ expansion of the meridional metric.
Numerical power-law fits yield empirical exponents that depend on the mode
set: $\hat\alpha=2.66$ for three modes, $\hat\alpha=2.37$ for five modes,
and $\hat\alpha=2.16$ for seven modes (Table~\ref{tab:power_law}).  The
empirical exponents converge towards the theoretical value as the mode count
increases, consistent with the expectation that higher-order modes more
effectively sample the curvature gradient and reduce finite-set artefacts.
The fitted prefactor is $\hat C\approx24$--$26$ across the tested mode sets,
while the analytical prediction from the $\sigma_\mathrm{min}$ expansion gives
$C_\mathrm{analytical}\approx160$.  The discrepancy reflects the limited
accuracy of the leading-order asymptotic approximation at moderate
$\varepsilon$.  All power-law fits are computed over the aspect-ratio range
$\zeta = c/a \in [0.1,\,0.95]$, corresponding to eccentricities
$\varepsilon\in[0.312,\,0.995]$; the modest coefficients of determination
($R^2\approx0.66$--$0.75$, Table~\ref{tab:power_law}) reflect both the
approximate nature of the power-law description and the curvature-channel floor
rounding off the divergence at high $\varepsilon$ (low $\zeta$).

\begin{table}[t]
  \centering
  \caption{Power-law fit $\kappa\approx C\,\varepsilon^{-\alpha}$ for different
  mode counts $\mathcal{N}$, using modes $n=2,\ldots,\mathcal{N}+1$ in the
  Ritz model.  Fits are over the aspect-ratio range $\zeta\in[0.1,\,0.95]$
  ($\varepsilon\in[0.312,\,0.995]$).  The theoretical exponent is $\alpha=2$;
  the modest $R^2$ values reflect the curvature-channel floor rounding off the
  divergence.}
  \label{tab:power_law}
  \begin{tabular}{cccc}
    \toprule
    $\mathcal{N}$ & $\hat\alpha$ & $\hat C$ & $R^2$ \\
    \midrule
    3 & 2.66 & 25.8 & 0.66 \\
    5 & 2.37 & 24.4 & 0.66 \\
    7 & 2.16 & 23.5 & 0.75 \\
    \bottomrule
  \end{tabular}
\end{table}

\paragraph{Ritz approximation error and rank restoration.}
A natural concern is whether the $\sim\!8\%$ approximation error of the Ritz
model relative to exact shell theory could itself be what breaks the
degeneracy, rather than the physical eccentricity.  Table~\ref{tab:mode_count}
addresses this indirectly.  If the rank restoration were a numerical artefact of
the Ritz truncation, one would expect the condition number to deteriorate (grow
towards singularity) as additional measured modes expose the approximation
error.  Instead, $\kappa$ remains stable across measured mode counts from
$\mathcal{N}=3$ to $\mathcal{N}=9$, ranging between $69$ and $94$, with no
trend towards divergence.  We therefore conclude that the identifiability
improvement is a genuine consequence of geometric eccentricity, not of
approximation error in the forward model.

\begin{figure}[t]
  \centering
  \includegraphics[width=0.85\linewidth]{fig_kappa_vs_eccentricity}
  \caption{Condition number $\kappa(\mathbf{J}_s)$ as a function of
  eccentricity $\varepsilon = \sqrt{1 - c^2/a^2}$, computed within the
  \emph{Ritz} model for five flexural modes ($n=2,\ldots,6$).  The condition
  number rises steeply as $\varepsilon \to 0$ (spherical limit), reaching a
  near-spherical floor of approximately $269$ (evaluated at
  $\varepsilon = 0.01$, i.e.\ $\zeta = 0.99995$, with five modes), and is
  approximately described by a power law
  $\kappa \approx C\,\varepsilon^{-\alpha}$ in the intermediate regime
  (fitted over $\zeta \in [0.1,\,0.95]$, i.e.\
  $\varepsilon\in[0.312,\,0.995]$) with theoretical exponent $\alpha=2$
  (empirical fit: $\hat\alpha=2.37$, $R^2=0.66$).  The green square marks the
  canonical operating point ($\zeta = 0.667$).  The result confirms the
  analytical predictions of Propositions~\ref{prop:chain_rule}
  and~\ref{prop:eccentricity}: oblate geometry restores identifiability through
  mode-dependent curvature sampling.}
  \label{fig:kappa_vs_eccentricity}
\end{figure}
